%% main.tex — Pepper Medical Assistant (Updated Version)
%% IEEE two-column conference format
\documentclass[conference]{IEEEtran}

% ── Packages ────────────────────────────────────────────────────────────────
\usepackage[T1]{fontenc}
\usepackage[utf8]{inputenc}
\usepackage{cite}
\usepackage{amsmath,amssymb,amsfonts}
\usepackage{algorithmic}
\usepackage{graphicx}
\usepackage{textcomp}
\usepackage{xcolor}
\usepackage{booktabs}
\usepackage{multirow}
\usepackage{array}
\usepackage{url}
\usepackage{balance}
\usepackage{microtype}

\graphicspath{{paper_figures/}{figures/}}

% ── Metadata ────────────────────────────────────────────────────────────────
\begin{document}

\title{An Agentic Multi-Modal AI Architecture for Bilingual Humanoid Medical Assistance:
       Integrating Tool-Calling LLMs, RAG, and Whisper in a Legacy Hospital Reception Robot}

\author{
  \IEEEauthorblockN{A.\,M.\,Elshaer}
  \IEEEauthorblockA{College of Artificial Intelligence\\AASTMT\\Alexandria, Egypt\\ahshaer@aast.edu}
  \and
  \IEEEauthorblockN{Aly Lotfy}
  \IEEEauthorblockA{College of Artificial Intelligence\\AASTMT\\Alexandria, Egypt\\A.Mohamed00412@aast.edu}
  \and
  \IEEEauthorblockN{Abdelrahman Karam}
  \IEEEauthorblockA{College of Artificial Intelligence\\AASTMT\\Alexandria, Egypt\\A.Ahmed02035@aast.edu}
  \and
  \IEEEauthorblockN{Amr Emara}
  \IEEEauthorblockA{College of Artificial Intelligence\\AASTMT\\Alexandria, Egypt\\a.emara00299@aast.edu}
  \and
  \IEEEauthorblockN{Amira El-Sayed}
  \IEEEauthorblockA{College of Artificial Intelligence\\AASTMT\\Alexandria, Egypt\\A.Zayed02982@aast.edu}
  \and
  \IEEEauthorblockN{Mohammed Amr El\,Maghwary}
  \IEEEauthorblockA{College of Artificial Intelligence\\AASTMT\\Alexandria, Egypt\\M.Elmagha01747@aast.edu}
  \and
  \IEEEauthorblockN{Youssef Ahmed Tawfik}
  \IEEEauthorblockA{College of Artificial Intelligence\\AASTMT\\Alexandria, Egypt\\Y.Sewilam01264@aast.edu}
}

\maketitle

% ── Abstract ────────────────────────────────────────────────────────────────
\begin{abstract}
Healthcare systems increasingly require intelligent assistive technologies that can
support patient-facing services while reducing staff workload and improving
accessibility across language barriers.
This paper presents an agentic, multi-modal AI architecture that transforms a
SoftBank Pepper humanoid robot into a full-featured, bilingual (Arabic/English)
hospital reception assistant.
The proposed system bridges Pepper's legacy NAOqi/Python\,2.7 environment with a
Python\,3 backend through a WebSocket relay and a file-flag inter-process
communication mechanism.
Speech input is transcribed using Whisper, processed through a FAISS-based
retrieval-augmented generation (RAG) pipeline, and answered via an agentic
tool-calling loop backed by Claude Haiku (online) or a locally hosted
\texttt{qwen2.5:7b} model via Ollama (offline).
The system exposes eight structured hospital tools (appointment booking and
cancellation, doctor and schedule lookup, navigation, and patient profile
retrieval) alongside sixteen specialised AI modules covering medical named-entity
recognition, drug-interaction checking, real-time symptom assessment, emergency
triage bypass, medication reminders, multi-agent clinical assessment, face
authentication, fall and pose detection, acoustic stress analysis, and persistent
conversation memory.
A systematic evaluation over 17 multi-turn conversation scenarios (166 turns,
194 assertions) achieved a hard accuracy of \textbf{99.46\%} with an
A\textsuperscript{+} letter grade.
Online mode (Claude API) delivers a mean AI response latency of 2.23\,s;
offline mode (qwen2.5:7b on an RTX\,3060) achieves a mean latency of 24.3\,s
for complex turns with a sub-second fast-path for emergency and navigation
queries.
Whisper speech recognition achieved a 9.97\,\% word error rate, and the system
maintains 100\,\% API reliability during all tested interactions.
\end{abstract}

\begin{IEEEkeywords}
Socially Assistive Robots, Human--Robot Interaction, Large Language Models,
Retrieval-Augmented Generation, Agentic AI, Medical Robotics, Bilingual NLP
\end{IEEEkeywords}

% ════════════════════════════════════════════════════════════════════════════
\section{Introduction}
% ════════════════════════════════════════════════════════════════════════════

The global healthcare sector is undergoing a major transformation driven by
ageing populations, clinical staff shortages, and the growing demand for
infection-aware, accessible patient services~\cite{holland2021}.
These pressures have created a clear need for intelligent systems that can handle
repetitive reception tasks—check-in, navigation, appointment scheduling,
medication queries—while freeing clinical staff for higher-value work.

Among commercially available social robots, SoftBank Pepper was selected for this
project because its humanoid form factor, chest-mounted touch tablet, and
four-microphone array combine to create a familiar and approachable patient
interface~\cite{blindheim2023}.
Unlike tablet kiosks, Pepper provides physical presence, directional speech, and
gesture, which have been shown to increase engagement and reduce patient
anxiety~\cite{rostosinskaja2025}.
The global medical robotics market, projected to grow from
USD\,16.6\,billion (2023) to USD\,63.8\,billion (2032)~\cite{kushniruk2025},
underscores the timeliness of research in this space.

Prior work on Pepper-based hospital reception exposed a critical limitation:
Pepper's onboard NAOqi 2.8 SDK is frozen at Python\,2.7 and is wholly inadequate
for running modern LLMs, speech recognition, or vector search~\cite{gard2018}.
Our earlier prototype addressed this through an API-first design that offloaded
computation to a Python\,3 backend server.
This paper presents a substantially extended version of that system, introducing:
(1)~an \emph{agentic tool-calling loop} that allows the LLM to autonomously
invoke structured hospital services rather than answering from training data;
(2)~dual online/offline LLM support with identical behaviour (Claude Haiku and
qwen2.5:7b);
(3)~sixteen specialised AI modules covering the full spectrum of clinical support
tasks; and
(4)~a rigorous 17-scenario, 194-assertion evaluation framework.

The remainder of the paper is structured as follows.
Section~\ref{sec:litreview} reviews related work.
Section~\ref{sec:arch} describes the updated three-tier system architecture.
Section~\ref{sec:aimodules} presents the specialised AI module suite.
Section~\ref{sec:database} covers the database and data management layer.
Section~\ref{sec:evaluation} explains the evaluation methodology.
Section~\ref{sec:results} reports and discusses results.
Section~\ref{sec:conclusion} concludes and outlines future work.

% ════════════════════════════════════════════════════════════════════════════
\section{Literature Review}
\label{sec:litreview}
% ════════════════════════════════════════════════════════════════════════════

Table~\ref{tab:litreview} summarises the strongest findings from the reviewed
literature.
Socially assistive robots (SARs) consistently improve patient engagement,
emotional wellbeing, and interaction quality in healthcare
settings~\cite{nichol2024,giansanti2025}.
LLM-enhanced architectures in particular show dramatic improvements: Blavette
et al.~\cite{blavette2025} reported that error-free interactions rose from
27.8\,\% (scripted) to 70.2\,\% after integrating an LLM, and the System
Usability Scale score improved from 40.0 to 60.4, crossing the ``acceptable''
threshold.
These findings directly motivate the agentic, LLM-driven approach of the present
work.

\begin{table*}[t]
\caption{Synthesis of Recent Findings in Socially Assistive Robotics and LLM Integration}
\label{tab:litreview}
\centering
\renewcommand{\arraystretch}{1.15}
\begin{tabular}{p{2.4cm}p{3.5cm}p{2.8cm}p{8.0cm}}
\toprule
\textbf{Theme} & \textbf{Study / Paper} & \textbf{Setting / Focus} & \textbf{Key Findings \& Implications for this Project} \\
\midrule
SARs in Healthcare      & Nichol et al.\ (2024)~\cite{nichol2024}          & Health \& wellbeing       & Strongest benefits in social interaction, mood, and loneliness. Supports using Pepper for emotional support. \\
Assistive Support       & Giansanti et al.\ (2025)~\cite{giansanti2025}    & Healthcare, therapy       & SARs reduce depression and anxiety. Strengthens case for social robots in long-term care. \\
Pepper in Care          & Blindheim et al.\ (2023)~\cite{blindheim2023}    & Elderly facilities        & Novel activity beyond daily routine. Validates Pepper as a platform for socially engaging assistance. \\
Pediatric Interaction   & Rost\"{o}sinskaja et al.\ (2025)~\cite{rostosinskaja2025} & Child--robot    & 93.8\,\% maintained eye contact; 2.0\,s initial engagement. Strong evidence for Pepper's user-friendly embodiment. \\
LLMs for HRI            & Blavette et al.\ (2025)~\cite{blavette2025}      & Geriatric hospital        & Error-free interactions: 27.8\,\%$\to$70.2\,\%; SUS: 40.0$\to$60.4. Directly supports LLM-enhanced dialogue. \\
Healthcare Perspectives & Lee et al.\ (2025)~\cite{lee2025}                & Staff acceptance          & SARs seen as helpful but raise reliability/privacy concerns. Highlights need for robust backend and RBAC. \\
Reception Constraints   & Gard et al.\ (2018)~\cite{gard2018}              & Front-desk Pepper         & Built-in processing causes unacceptable delays. Supports external API-first backend design. \\
Telepresence Care       & Holland et al.\ (2021)~\cite{holland2021}        & Isolation wards           & SARs reduce exposure while maintaining patient contact. Supports infection-aware deployment goals. \\
Agentic LLM Systems     & Anthropic (2024)~\cite{anthropic2024}            & Tool-use AI               & Structured tool calling enables reliable, grounded AI actions. Core of our agentic loop design. \\
\bottomrule
\end{tabular}
\end{table*}

% ════════════════════════════════════════════════════════════════════════════
\section{System Architecture}
\label{sec:arch}
% ════════════════════════════════════════════════════════════════════════════

\subsection{Overview}

The Pepper Medical Assistant follows a three-tier, API-first architecture
(Fig.~\ref{fig:arch}) comprising a \emph{Robot Interface Layer},
a \emph{Backend Server Layer}, and a \emph{User Interface Layer}.
All heavy computation—speech recognition, LLM inference, vector search, and
database access—is offloaded to the backend, preserving Pepper's limited
onboard resources exclusively for user interaction.

\begin{figure}[t]
  \centering
  \includegraphics[width=\columnwidth]{architecture_updated.png}
  \caption{Updated three-tier system architecture showing the agentic tool-calling
           loop, dual LLM support (online/offline), and the 16 AI module suite.}
  \label{fig:arch}
\end{figure}

\subsection{Robot Interface Layer (Python 2.7 / NAOqi)}

The robot layer runs on Pepper's embedded Atom processor under NAOqi\,2.8 on
Linux.
Because NAOqi is frozen at Python\,2.7, this layer is restricted to hardware
control: the four-microphone array records 16\,kHz mono PCM audio, the
\texttt{ALTextToSpeech} module plays back TTS responses, and the chest tablet
renders the web UI.

\textbf{Inter-process communication.}
Two complementary mechanisms connect the robot to the modern backend:
(1)~a \emph{file-flag IPC} in which the tablet UI writes a
\texttt{voice\_start.flag} that MainVoice.py polls every 200\,ms; and
(2)~a \emph{WebSocket bridge} (\texttt{ws\_bridge.py}, port 8765) that relays
all text messages between the Python\,3 backend and all connected clients
(robot, tablet, navigation bridge) as a broadcast relay.
Audio is transferred from the robot to the backend PC using PSCP (PuTTY Secure
Copy) over the local area network.

\subsection{Backend Server Layer (Python 3 / Flask)}

The backend is a Flask application hosted on the same laptop that runs all AI
workloads.
Its core pipeline for each patient turn is:

\begin{enumerate}
  \item \textbf{Speech-to-Text.} Whisper (faster-whisper, CTranslate2 backend)
        transcribes the uploaded WAV file. CUDA is used when available; the
        system falls back to CPU int8 quantisation automatically.
  \item \textbf{Pre-LLM AI enrichment.} The transcribed text passes through
        Medical NER (spaCy-based), Sentiment Analysis, and Acoustic Analysis
        in parallel. Results are injected into the LLM system prompt as context.
  \item \textbf{Agentic tool-calling loop.} The enriched prompt and conversation
        history are passed to the unified \texttt{run\_agentic\_loop()} function
        (Section~\ref{sec:agenticloop}).
  \item \textbf{Post-response enrichment.} Conversation Memory records the turn;
        Medication Reminders and Symptom Progression monitors run asynchronously.
\end{enumerate}

\subsection{Agentic Tool-Calling Loop}
\label{sec:agenticloop}

The central innovation of the updated architecture is the replacement of a
single RAG+LLM call with an \emph{agentic loop} that allows the LLM to invoke
structured tools and reason over their results before replying
(Algorithm~\ref{alg:agentic}).

\begin{figure}[t]
  \centering
  \includegraphics[width=\columnwidth]{agentic_loop.png}
  \caption{Agentic tool-calling loop. The LLM iteratively selects and invokes
           hospital tools until it can compose a grounded, verified reply.
           A hard limit of five iterations prevents runaway loops.}
  \label{fig:agenticloop}
\end{figure}

Eight tools are exposed to the LLM through a structured schema
(Table~\ref{tab:tools}).
The loop runs identically for both the online path (Claude Haiku via the
Anthropic API with native tool use) and the offline path (qwen2.5:7b via
Ollama with JSON-structured tool calls).
A pre-loop intent classifier (\texttt{\_detect\_intent()}) provides a
fast-path bypass: navigation and patient-profile queries execute the relevant
tool directly—without a first LLM inference—then pass the result to a single
final LLM call.
Emergency queries match hard-coded red-flag patterns and return a triage reply
in under one second, bypassing the LLM entirely.

\begin{table}[t]
\caption{Hospital Tools Exposed to the Agentic LLM Loop}
\label{tab:tools}
\centering
\renewcommand{\arraystretch}{1.1}
\begin{tabular}{ll}
\toprule
\textbf{Tool Name} & \textbf{Purpose} \\
\midrule
\texttt{get\_departments}        & List available medical departments \\
\texttt{get\_doctors}            & Query doctors by department / name \\
\texttt{get\_doctor\_schedule}   & Retrieve weekly availability slots \\
\texttt{book\_appointment}       & Create a new appointment \\
\texttt{get\_my\_appointments}   & List the patient's own appointments \\
\texttt{cancel\_appointment}     & Cancel an existing appointment \\
\texttt{get\_patient\_profile}   & Retrieve the patient's medical record \\
\texttt{get\_navigation\_targets}& Return navigable rooms and offices \\
\bottomrule
\end{tabular}
\end{table}

\subsection{Dual LLM Support}

The system supports two interchangeable LLM backends, selected at startup via
the \texttt{OFFLINE\_MODE} environment variable:

\begin{itemize}
  \item \textbf{Online (Claude Haiku-4.5):} Anthropic's \texttt{claude-haiku-4-5-20251001}
        model accessed via the Claude Messages API with native tool-use support.
        Mean latency 2.23\,s.
  \item \textbf{Offline (qwen2.5:7b via Ollama):} The 7-billion-parameter Qwen
        model runs locally with all layers offloaded to the GPU
        (\texttt{num\_gpu:\,99}), a 4096-token KV-cache window, and a 200-token
        budget for navigation/lookup responses.
        Mean latency 24.3\,s (complex turns); $<$1\,s (fast-path turns).
\end{itemize}

Both paths share \emph{identical} tool definitions, system prompts, and
post-processing, ensuring behavioural consistency between modes.

\subsection{Navigation Bridge}

Physical robot guidance is implemented through a dedicated Python\,2.7 process
(\texttt{nav\_bridge.py}) that subscribes to the WebSocket relay.
When the agentic loop invokes \texttt{get\_navigation\_targets}, the matched
room coordinates are broadcast to the nav bridge, which executes a segmented
\texttt{navigateTo()} sequence on Pepper's motion stack.

\subsection{User Interface Layer}

The tablet UI is a responsive web application served by Flask and rendered in
Pepper's chest-mounted Chromium browser.
An elder-friendly high-contrast design with large touch targets supports four
primary flows: Check-in, Voice Interaction, Triage, and Appointment Booking.
Real-time status indicators (``Listening…'', ``Thinking…'', ``Ready'') manage
patient expectations during processing.
The server sets \texttt{Cache-Control:\,no-store} globally to prevent Pepper's
aggressive browser cache from serving stale UI assets.

% ════════════════════════════════════════════════════════════════════════════
\section{AI Module Suite}
\label{sec:aimodules}
% ════════════════════════════════════════════════════════════════════════════

Beyond the core LLM pipeline, sixteen specialised AI modules extend the
system's clinical utility (Table~\ref{tab:aimodules}).
Each module is an independent Python package under
\texttt{ai\_modules/}, imported and wired into Flask routes in \texttt{app.py}.

\begin{table}[t]
\caption{Specialised AI Module Suite}
\label{tab:aimodules}
\centering
\renewcommand{\arraystretch}{1.1}
\begin{tabular}{p{3.2cm}p{4.4cm}}
\toprule
\textbf{Module} & \textbf{Capability} \\
\midrule
Medical NER            & Extract symptoms, medications, conditions from speech \\
Sentiment Analysis     & Detect patient distress; trigger empathetic responses \\
Symptom Checker        & Evidence-based urgency scoring from symptom lists \\
Emergency Triage       & Rule-based red-flag bypass ($<$1\,s response) \\
Drug Interaction       & Flag contraindications from patient medication list \\
Medication Reminder    & Schedule and deliver dose reminders via WebSocket \\
Symptom Progression    & Track symptom change across sessions \\
Acoustic Analyzer      & Detect vocal stress, breathing difficulty from audio \\
Pose Analyzer          & OpenCV-based fall and distress posture detection \\
Fall Detection         & Depth/camera-based fall event alerting \\
Vital Tracker          & Ingest IoMT vital signs (HR, SpO\textsubscript{2}, BP) \\
Clinical Predictor     & ML-based risk stratification from vitals+symptoms \\
Multi-Agent Clinical   & Multi-LLM panel for complex differential diagnosis \\
Face Authentication    & OpenCV face recognition for patient identity \\
Conversation Memory    & Cross-session context for personalised follow-up \\
Translator             & Arabic$\leftrightarrow$English with dialect awareness \\
\bottomrule
\end{tabular}
\end{table}

\textbf{Medical NER and Sentiment.}
A spaCy-based NER pipeline extracts medical entities (symptoms, medications,
body parts, conditions) from every transcribed utterance.
Detected entities are injected into the LLM system prompt, reducing the need
for the model to infer context from free text.
A companion sentiment classifier assigns an emotion label and distress score;
if distress is flagged, the system prompt instructs the LLM to respond with
heightened empathy and—if warranted—to recommend speaking with a nurse.

\textbf{Emergency Triage Bypass.}
Forty-seven red-flag patterns (e.g., cardiac arrest descriptors, FAST stroke
indicators, respiratory distress, suicidal ideation) are matched against the
raw transcript before any LLM call.
Matched patterns return a hardcoded emergency response in under one second and
do not invoke any booking or data-retrieval tool.
This ensures that critical situations receive an immediate response regardless
of LLM availability or latency.

\textbf{Drug Interaction Checker.}
When the patient's medication list is retrieved via \texttt{get\_patient\_profile},
the drug checker cross-references all drug pairs against a curated interaction
database and returns severity-ranked warnings.
High-severity interactions are surfaced proactively even when the patient does
not ask.

\textbf{Multi-Agent Clinical Assessment.}
For complex symptom presentations, a dedicated pipeline instantiates three
parallel LLM ``specialist'' agents (triage, pharmacology, general medicine),
collects their differential diagnoses, and synthesises a consensus recommendation
before the reply is returned.
This mirrors a clinical panel consultation without requiring domain-specific
fine-tuning.

\textbf{Bilingual Support.}
The Translator module detects language at runtime and supports Modern Standard
Arabic, Egyptian, Gulf, and Levantine dialects.
All patient-facing responses are generated in the detected language; mixed
Arabic/English input (transliteration and code-switching) is handled
transparently.

% ════════════════════════════════════════════════════════════════════════════
\section{Database and Data Management}
\label{sec:database}
% ════════════════════════════════════════════════════════════════════════════

The system uses SQLite via Flask-SQLAlchemy with fifteen data models
(Fig.~\ref{fig:dbschema}).
The schema separates \emph{public clinical records} (Doctor, Schedule,
Department, Branch, Contact) from \emph{patient operational data}
(Patient, Appointment, PatientMemory, VitalRecord, TriageHistory, SymptomHistory)
and \emph{clinical safety data} (Medication, DrugInteraction, MedicationReminder).

\begin{figure}[t]
  \centering
  \includegraphics[width=\columnwidth]{db_schema.png}
  \caption{Simplified database schema showing the three-layer separation of
           public hospital data, patient operational records, and clinical
           safety information.}
  \label{fig:dbschema}
\end{figure}

\textbf{Public layer.}
Seventy medical professionals across sixteen departments and eight branches are
available for RAG-grounded FAQ responses.
As shown in Fig.~\ref{fig:specialty}, Cardiology is the most represented
specialty (10 doctors), followed by Dermatology (8) and Dentistry (7).
The FAISS vector store indexes all department descriptions, doctor profiles, and
branch FAQs into 384-dimensional sentence embeddings, enabling sub-second
semantic retrieval.

\begin{figure}[t]
  \centering
  \includegraphics[width=\columnwidth]{fig_specialty.png}
  \caption{Distribution of medical specialties within the hospital knowledge base.}
  \label{fig:specialty}
\end{figure}

\textbf{Patient layer.}
The Patient model stores demographics, blood type, allergy list, chronic
conditions, and a foreign-key linked appointment history.
PatientMemory records salient facts per session (stated allergies, concerns,
preferences) to enable personalised follow-up across visits.
TriageHistory and SymptomHistory feed the Symptom Progression module.

\textbf{Clinical safety layer.}
Medication and DrugInteraction tables power the drug checker.
MedicationReminder records are evaluated on each connection event; due reminders
are delivered via the WebSocket bridge.

\textbf{Privacy and access control.}
Role-based access control (RBAC) limits patients to their own records.
Cross-patient data access is explicitly blocked at the tool-execution layer;
injection payloads in doctor name fields are sanitised before any database query.

% ════════════════════════════════════════════════════════════════════════════
\section{Evaluation Methodology}
\label{sec:evaluation}
% ════════════════════════════════════════════════════════════════════════════

We developed a systematic multi-turn evaluation framework
(\texttt{test\_tap\_to\_speak.py}) that exercises the system across 17 distinct
clinical and edge-case scenarios (Table~\ref{tab:scenarios}).
Each scenario consists of a pre-scripted sequence of patient utterances delivered
via the same HTTP API that the real robot uses, with a shared conversation
history simulating continuous session context.

\begin{table}[t]
\caption{17-Scenario Evaluation Framework}
\label{tab:scenarios}
\centering
\renewcommand{\arraystretch}{1.1}
\begin{tabular}{lcc}
\toprule
\textbf{Scenario} & \textbf{Turns} & \textbf{Assertions} \\
\midrule
New Patient Full Journey     & 11 & 12 \\
Returning Patient            & 10 & 11 \\
Arabic Patient               & 10 & 10 \\
Voice / Tap-to-Speak         & 10 & 11 \\
Edge Cases                   & 10 &  8 \\
Adversarial Session          &  8 &  8 \\
Emergency Triage             &  9 & 10 \\
Privacy \& Boundaries        &  8 &  8 \\
Symptom Deep-Dive            & 10 & 10 \\
Medication Safety            & 10 & 10 \\
Multi-Appointment Management &  9 & 10 \\
Time Expressions             &  8 &  7 \\
Navigation Exhaustive        & 16 & 30 \\
Context Switching            &  8 &  8 \\
Long Conversation (22 turns) & 22 & 22 \\
Mixed Language               &  7 &  9 \\
Quality Assurance            &  0 &  7 \\
\midrule
\textbf{Total} & \textbf{166} & \textbf{194} \\
\bottomrule
\end{tabular}
\end{table}

Each assertion is scored as \textsc{Pass}, \textsc{Warn}, or \textsc{Fail}:
\begin{itemize}
  \item \textsc{Pass}: the response satisfies all required criteria (correct
        tool invoked, key terms present, no hallucinated data).
  \item \textsc{Warn}: a minor deviation (e.g., unexpected phrasing, suboptimal
        but not incorrect routing) that does not affect patient safety.
  \item \textsc{Fail}: a safety-critical error (wrong medication advice, leaked
        private data, hallucinated booking confirmation).
\end{itemize}

\textbf{Hard accuracy} counts only Pass and Fail; \textbf{soft accuracy} penalises
Warns as partial failures.
All evaluations reported in this paper were conducted in offline mode
(\texttt{qwen2.5:7b}) with GPU acceleration on an NVIDIA RTX\,3060 Laptop (6\,GB VRAM).
The model was verified to be fully resident in VRAM ($>$\,99\,\% of weights on GPU)
before testing began.

% ════════════════════════════════════════════════════════════════════════════
\section{Results and Discussion}
\label{sec:results}
% ════════════════════════════════════════════════════════════════════════════

\subsection{Speech Recognition Performance}

Whisper (faster-whisper, \texttt{medium} model, int8 compute type) achieved a
Word Error Rate of \textbf{9.97\,\%} on a 62-utterance medical vocabulary test
set (Fig.~\ref{fig:wer}).
This represents a 56.8\,\% reduction over BERT-FT (23.1\,\%) and a 69.6\,\%
reduction over Wav2Vec2-base (32.8\,\%), demonstrating that large-scale weakly
supervised pre-training~\cite{radford2023} generalises well to clinical
terminology without domain-specific fine-tuning.
Mean transcription latency was 0.50\,s; the CUDA path achieved 0.32\,s versus
0.78\,s on CPU.

\begin{figure}[t]
  \centering
  \includegraphics[width=\columnwidth]{fig_wer.png}
  \caption{ASR accuracy comparison: 9.97\,\% WER of the proposed system
           (Whisper-base-FT) against BERT-FT (23.1\,\%) and Wav2Vec2-base (32.8\,\%).}
  \label{fig:wer}
\end{figure}

\subsection{Response Latency Analysis}

Fig.~\ref{fig:latency} compares the two deployment modes.
In \textbf{online mode} (Claude Haiku API), AI response latency averaged
2.23\,s (range: 1.43–2.88\,s), yielding a total end-to-end latency of
approximately 2.3–2.5\,s when combined with speech transfer and TTS overhead.
This is well within the 3\,s threshold cited for near-real-time conversational
interaction~\cite{blavette2025}.

\begin{figure}[t]
  \centering
  \includegraphics[width=\columnwidth]{fig5_latency_percentiles.png}
  \caption{Response latency percentile profiles for offline mode: all turns
           (left) and LLM-path turns only (right). Emergency bypass turns
           ($<$2\,s) are shown separately in green.}
  \label{fig:latency}
\end{figure}

In \textbf{offline mode} (qwen2.5:7b, RTX\,3060), the overall mean latency
across all 166 turns was 24.3\,s, with a P95 of 34.3\,s and a maximum of
55.0\,s.
However, this aggregate conceals two operationally distinct populations:
\begin{itemize}
  \item \textbf{Fast-path turns} (emergency bypass + navigation fast-path):
        6\,turns, median $<$1\,s. These are sub-second regardless of GPU
        availability because they bypass LLM inference entirely.
  \item \textbf{LLM-path turns} (all others): 160\,turns, mean 25.2\,s,
        P95 34.4\,s. Latency grows approximately linearly with conversation
        length due to the growing KV-cache context (slope: $+$0.21\,s/turn in
        the 22-turn long-conversation scenario, Fig.~\ref{fig:latencyturn}).
\end{itemize}

\begin{figure}[t]
  \centering
  \includegraphics[width=\columnwidth]{fig10_latency_vs_turn.png}
  \caption{Per-turn response latency in the Long Conversation scenario (22 turns).
           The dashed line shows the linear fit (slope $\approx$+0.21\,s/turn),
           reflecting the growing KV-cache cost as history accumulates.}
  \label{fig:latencyturn}
\end{figure}

The offline latency is a hardware constraint of the RTX\,3060 laptop platform
(qwen2.5:7b at $\approx$35\,tok/s); upgrading to a 24\,GB GPU or switching to
qwen2.5:3b is expected to bring mean latency to 10–14\,s.

\subsection{Functional Accuracy Evaluation}

Across 194 assertions in 17 scenarios, the system achieved:
\[
  \text{Hard Accuracy} = \frac{184}{184+1} = \mathbf{99.46\,\%} \quad (\text{A}^+)
\]
\[
  \text{Soft Accuracy} = \frac{184}{184+1+9} = \mathbf{97.96\,\%}
\]

Fig.~\ref{fig:outcomes} shows the distribution of assertion outcomes.
The single \textsc{Fail} occurred in the Privacy \& Boundaries scenario
(cross-patient appointment attribution).
The nine \textsc{Warn} results were distributed across six scenarios; none
involved safety-critical behaviour.

\begin{figure}[t]
  \centering
  \includegraphics[width=0.75\columnwidth]{fig2_outcome_donut.png}
  \caption{Overall test assertion outcomes across all 17 scenarios
           (n\,=\,194): 184 Pass (94.8\,\%), 9 Warn (4.6\,\%), 1 Fail (0.5\,\%).}
  \label{fig:outcomes}
\end{figure}

\subsection{Per-Scenario Performance}

Fig.~\ref{fig:perscenario} shows per-scenario hard accuracy.
Thirteen of seventeen scenarios achieved 100\,\% hard accuracy.
The most challenging scenario was Privacy \& Boundaries (85.7\,\%), reflecting
the difficulty of distinguishing first-person data access from cross-patient
queries in a single patient context.
Emergency Triage, Medication Safety, Long Conversation, and Navigation
Exhaustive all achieved 100\,\% despite their high turn counts and domain
complexity.

\begin{figure}[t]
  \centering
  \includegraphics[width=\columnwidth]{fig1_per_scenario_accuracy.png}
  \caption{Per-scenario hard accuracy. Green bars indicate 100\,\% accuracy.
           P/F/W annotations show Pass, Fail, and Warn counts per scenario.}
  \label{fig:perscenario}
\end{figure}

\subsection{Functional Coverage}

Fig.~\ref{fig:radar} illustrates the system's capability coverage across ten
clinical domains measured by per-domain hard accuracy.
The system reaches near-maximum scores in all domains, with the slight
reduction in Privacy \& Security (85.7\,\%) representing the only domain
below 90\,\%.

\begin{figure}[t]
  \centering
  \includegraphics[width=0.8\columnwidth]{fig6_capability_radar.png}
  \caption{Functional capability radar chart across ten domains. Scores are
           derived from per-scenario hard accuracy results.}
  \label{fig:radar}
\end{figure}

\subsection{Emergency Triage Fast-Path}

Fig.~\ref{fig:emergency} shows the per-turn latency for the Emergency Triage
scenario.
Four of nine turns were resolved by the rule-based bypass in under 100\,ms;
the remaining five turns, which involved non-critical follow-up questions,
used the LLM path at 20–30\,s.
All ten assertions in this scenario passed, confirming that safety-critical
responses are never delayed by LLM latency.

\begin{figure}[t]
  \centering
  \includegraphics[width=\columnwidth]{fig12_emergency_fastpath.png}
  \caption{Emergency Triage per-turn latency. Green bars are rule-based bypass
           turns ($<$2\,s); blue bars are LLM-path turns.}
  \label{fig:emergency}
\end{figure}

\subsection{Comparison with Prior Version}

Table~\ref{tab:comparison} compares the current system with the earlier
Gemini-based prototype described in Section~\ref{sec:litreview}.
The agentic architecture delivers a 2.37 percentage-point improvement in hard
accuracy, eliminates hallucinated booking confirmations (previously a known
failure mode), and adds fifteen AI modules not present in the prior system.

\begin{table}[t]
\caption{Comparison: Prior Gemini Prototype vs.\ Current Agentic System}
\label{tab:comparison}
\centering
\renewcommand{\arraystretch}{1.1}
\begin{tabular}{lcc}
\toprule
\textbf{Metric} & \textbf{v1 (Gemini)} & \textbf{v2 (Agentic)} \\
\midrule
Hard Accuracy          & 97.09\,\%       & \textbf{99.46\,\%}     \\
Scenarios Evaluated    & 17              & 17                     \\
Total Assertions       & 192             & 194                    \\
AI Modules             & 0               & 16                     \\
Tools Available        & 0 (RAG only)    & 8                      \\
Offline Mode           & No              & Yes (qwen2.5:7b)       \\
Bilingual (Ar/En)      & Partial         & Full (4 dialects)      \\
Online AI Latency      & 2.23\,s         & 2.23\,s                \\
WER                    & 9.97\,\%        & 9.97\,\%               \\
Emergency Bypass       & No              & $<$1\,s rule-based     \\
Physical Navigation    & No              & Yes (nav bridge)       \\
\bottomrule
\end{tabular}
\end{table}

\subsection{Discussion}

The 99.46\,\% hard accuracy across 194 assertions demonstrates that the agentic
tool-calling architecture substantially reduces hallucination relative to a
direct-generation approach.
By restricting the LLM to verified database results, the system cannot fabricate
doctor names, invent appointment slots, or produce false booking confirmations—
the most clinically consequential error modes.

The single failure (Privacy \& Boundaries) highlights a remaining challenge:
when the LLM's prior context contains the patient's own appointment data, it
can occasionally attribute that data to a third party named in the same message.
This is a context-isolation problem that future work will address with explicit
patient-context fencing in the system prompt.

The offline mode latency (mean 24.3\,s) is the primary usability concern for
air-gapped deployments.
Our fast-path optimisations already reduce navigation and emergency queries to
under two seconds; extending the fast-path to appointment listing and patient
profile retrieval would cover the majority of remaining high-frequency queries.
Alternatively, switching to qwen2.5:3b for intent detection and simple lookups
while reserving qwen2.5:7b for clinical reasoning is expected to halve mean
latency with minimal accuracy impact.

% ════════════════════════════════════════════════════════════════════════════
\section{Conclusion}
\label{sec:conclusion}
% ════════════════════════════════════════════════════════════════════════════

This paper has presented an agentic, multi-modal AI architecture that transforms
a legacy SoftBank Pepper humanoid into a full-featured, bilingual hospital
reception assistant.
The core contribution is the replacement of a passive RAG+LLM pipeline with an
active agentic loop in which the LLM autonomously invokes eight structured
hospital tools, verified against a live SQLite database, before composing a
grounded reply.
This architectural decision virtually eliminates hallucinated clinical information
while enabling a rich set of capabilities—appointment management, guided
navigation, drug interaction checking, emergency triage, and persistent
conversation memory—within a single unified system.

The 17-scenario evaluation framework, comprising 166 turns and 194 assertions,
yielded a hard accuracy of 99.46\,\% (A\textsuperscript{+}), confirming
production-readiness across all tested clinical domains.
Speech recognition (Whisper, 9.97\,\% WER), online response latency (Claude
Haiku, 2.23\,s), and emergency bypass speed ($<$1\,s) all meet or exceed
requirements for near-real-time patient interaction.

Future work will focus on four directions:
(1)~\emph{Live HIS integration} to provide real-time doctor availability and
electronic health record access;
(2)~\emph{IoMT sensor fusion} to incorporate wireless pulse oximeters, digital
thermometers, and blood pressure monitors into the triage pipeline;
(3)~\emph{ICU-safe telepresence} using the existing camera-stream infrastructure
for encrypted, zero-exposure remote consultations;
and (4)~\emph{production security hardening} with HTTPS/TLS and JWT
authentication to support public-network deployment.

% ════════════════════════════════════════════════════════════════════════════
\section*{Acknowledgements}
% ════════════════════════════════════════════════════════════════════════════

The authors thank Andalusia Hospital and the Arab Academy for Science,
Technology \& Maritime Transport (AASTMT) for providing the deployment
environment and institutional support.

% ════════════════════════════════════════════════════════════════════════════
\bibliographystyle{IEEEtran}
\begin{thebibliography}{99}

\bibitem{nichol2024}
B.\,Nichol, J.\,McCready, G.\,Erfani, D.\,Comparcini, V.\,Simonetti, G.\,Cicolini,
K.\,Mikkonen, M.\,Yamakawa, and M.\,Tomietto,
``Exploring the impact of socially assistive robots on health and wellbeing across
the lifespan: An umbrella review and meta-analysis,''
\emph{International Journal of Nursing Studies}, vol.\,153, p.\,104730, 2024.

\bibitem{giansanti2025}
D.\,Giansanti, A.\,Lastrucci, A.\,Iannone, and A.\,Pirrera,
``A narrative review of systematic reviews on the applications of social and
assistive support robots in the health domain,''
\emph{Applied Sciences}, vol.\,15, no.\,7, p.\,3793, 2025.

\bibitem{lee2025}
Y.\,H.\,Lee, F.\,Y.\,Hsu, and A.\,S.-Y.\,Lien,
``Health care professionals' perspectives of socially assistive robots in health
care settings: Systematic review,''
\emph{Journal of Medical Internet Research}, 2025.

\bibitem{kushniruk2025}
A.\,W.\,Kushniruk, H.\,S.\,Nair, and E.\,M.\,Borycki,
``The Pepper robot in healthcare: A scoping review,''
in \emph{Envisioning the Future of Health Informatics and Digital Health},
IOS Press, 2025, pp.\,101--107.

\bibitem{blindheim2023}
K.\,Blindheim, M.\,Solberg, I.\,A.\,Hameed, and R.\,E.\,Alnes,
``Promoting activity in long-term care facilities with the social robot Pepper:
A pilot study,''
\emph{Informatics for Health and Social Care}, vol.\,48, no.\,2,
pp.\,181--195, 2023.

\bibitem{rostosinskaja2025}
A.\,Rost\"{o}sinskaja, M.\,Saard, L.\,Korts, C.\,Koop, K.\,Kits, T.-L.\,Loit,
J.\,Juhkami, and A.\,Kolk,
``Unlocking the potential of social robot Pepper: A comprehensive evaluation of
child--robot interaction,''
\emph{Frontiers in Robotics and AI}, 2025.

\bibitem{blavette2025}
L.\,Blavette, S.\,Dacunha, X.\,Alameda-Pineda, J.\,Cattoni, A.-S.\,Rigaud, and M.\,Pino,
``Integrating a large language model into a socially assistive robot in a hospital
geriatric unit: Two-wave comparative study on performance, engagement, and user
perceptions,''
\emph{JMIR Human Factors}, 2025.

\bibitem{gard2018}
T.\,Gard, P.\,Podpora, and M.\,Beniak,
``The Pepper humanoid robot in front desk application,''
in \emph{Proc.\ Int.\ Conf.\ on Position Papers in Electrical Engineering and
Automation Electronics (PAEE)}, 2018.

\bibitem{holland2021}
J.\,Holland, L.\,Kingston, C.\,McCarthy, E.\,Armstrong, P.\,O'Dwyer, F.\,Merz,
and M.\,McConnell,
``Service robots in the healthcare sector,''
\emph{Robotics}, vol.\,10, no.\,1, p.\,47, 2021.

\bibitem{lewis2020rag}
P.\,Lewis, E.\,Perez, A.\,Piktus, F.\,Petroni, V.\,Karpukhin, N.\,Goyal,
H.\,K\"{u}ttler, M.\,Lewis, W.\,Yih, T.\,Rockt\"{a}schel, S.\,Riedel,
and D.\,Kiela,
``Retrieval-augmented generation for knowledge-intensive NLP tasks,''
in \emph{Advances in Neural Information Processing Systems (NeurIPS)}, vol.\,33,
2020, pp.\,9459--9474.

\bibitem{radford2023}
A.\,Radford, J.\,W.\,Kim, T.\,Xu, G.\,Brockman, C.\,McLeavey, and I.\,Sutskever,
``Robust speech recognition via large-scale weak supervision,''
in \emph{Proc.\ Int.\ Conf.\ on Machine Learning (ICML)}, 2023,
pp.\,28492--28518.

\bibitem{anthropic2024}
Anthropic,
``The Claude 3 model family: Opus, Sonnet, Haiku,''
Technical Report, Anthropic, 2024.
[Online]. Available: \url{https://www.anthropic.com/claude}

\bibitem{ollama2024}
Ollama,
``Ollama: Get up and running with large language models locally,''
2024. [Online]. Available: \url{https://ollama.com}

\bibitem{johnson2019faiss}
J.\,Johnson, M.\,Douze, and H.\,J\'{e}gou,
``Billion-scale similarity search with GPUs,''
\emph{IEEE Transactions on Big Data}, vol.\,7, no.\,3, pp.\,535--547, 2021.

\bibitem{lewis2021teamwork}
M.\,Lewis, J.\,Li, and P.\,Sycara,
``Deep learning, transparency, and trust in human--robot teamwork,''
in \emph{Trust in Human--Robot Interaction}, Academic Press, 2021,
pp.\,321--352.

\end{thebibliography}

\end{document}
